\subsection{DETERMINANT OF A SQUARE MATRIX}

DEFINITION 2. Let I be a finite set, A a commutative ring and X a square matrix of type (I, I) over the ring A (II, § 10, no. 7). The determinant of the endomorphism u of the A-module A \(^{I}\) , whose matrix with respect to the canonical basis of A \(^{I}\) is X, is called the determinant of X and denoted by det X.

If \(X = (\xi_{ij})_{(i,j) \in I \times I}\) and \((e_i)_{i \in I}\) is the canonical basis of \(A^I\) , the endomorphism \(u\) is then given by

\[
u (e _ {i}) = \sum_ {j \in J} \xi_ {j i} e _ {j}.
\]

When \(\mathbf{I} = [1, n] \subset \mathbf{N}\) , if we write \(x_{i} = u(e_{i})\) for \(i \in \mathbf{I}\) , the determinant of \(X\) is then defined in the relation

\[
x _ {1} \wedge x _ {2} \wedge \dots \wedge x _ {n} = (\det X) e _ {1} \wedge e _ {2} \wedge \dots \wedge e _ {n}\tag{6}
\]

in other words, det X is equal to the determinant \(\det(x_{1}, x_{2}, \ldots, x_{n})\) with respect to the canonical basis of \(A^{n}\) . Consequently:

PROPOSITION 4. For \(n\) vectors \(x_1, \ldots, x_n\) of \(A^n\) , let \(X(x_1, \ldots, x_n)\) denote the square matrix of order \(n\) whose \(i\) -th column is \(x_i\) for \(1 \leqslant i \leqslant n\) . Then the mapping

\[
(x _ {1}, \dots , x _ {n}) \mapsto \det (X (x _ {1}, \dots , x _ {n}))
\]

of \((\mathbf{A}^n)^n\) into \(\mathbf{A}\) is alternating and \(n\) -linear.

In particular, the determinant of a matrix two of whose columns are equal is zero. If a permutation is performed on the columns of a matrix, the determinant of the new matrix is equal to that of the old multiplied by \(\varepsilon_{\sigma}\) . If to one column of a matrix is added a scalar multiple of a column of a different index, the determinant of the new matrix is equal to that of the old.

More generally, let M be a free A-module of finite dimension n and let \((e_{i})_{i\in I}\) be a basis of M; for every automorphism u of M, if X is the matrix of u with respect to the basis \((e_{i})\) , then

\[
\det (u) = \det (X)\tag{7}
\]

as follows immediately from the definitions.

When \(\mathbf{I} = [1, n]\) , the determinant of \(X\) is also denoted by

\[
\det (\xi_ {i j}) _ {1 \leqslant i \leqslant n, 1 \leqslant j \leqslant n}
\]

or simply \(\operatorname{det}(\xi_{ij})\) if this causes no confusion, or also

\[
\left| \begin{array}{c c c c c} \xi_ {1 1} & \xi_ {1 2} & \ldots & \xi_ {1 n} \\ \xi_ {2 1} & \xi_ {2 2} & \ldots & \xi_ {2 n} \\ \cdot & \cdot & \cdot & \cdot & \cdot \\ \xi_ {n 1} & \xi_ {n 2} & \ldots & \xi_ {n n} \end{array} \right|
\]

When \(X = 1\) , the matrix \(X\) is called unimodular.

Examples. (1) The determinant of the empty matrix is equal to 1; the determinant of a square matrix of order 1 is equal to the unique element of this matrix. For a matrix of order 2

\[
\left( \begin{array}{c c} \xi_ {1 1} & \xi_ {1 2} \\ \xi_ {2 1} & \xi_ {2 2} \end{array} \right)
\]

then, in the above notation,

\[
x _ {1} \wedge x _ {2} = (\xi_ {1 1} e _ {1} + \xi_ {2 1} e _ {2}) \wedge (\xi_ {1 2} e _ {1} + \xi_ {2 2} e _ {2}) = \xi_ {1 1} \xi_ {2 2} e _ {1} ^ {\prime} \wedge e _ {2} + \xi_ {2 1} \xi_ {1 2} e _ {2} \wedge e _ {1}
\]

whence

\[
\left| \begin{array}{c c} \xi_ {1 1} & \xi_ {1 2} \\ \xi_ {2 1} & \xi_ {2 2} \end{array} \right| = \xi_ {1 1} \xi_ {2 2} - \xi_ {1 2} \xi_ {2 1}.
\]

We translate into the language of matrices some of the results of nos. 1 and 2:

PROPOSITION 5. If X and Y are two square matrices over a commutative ring A with the same finite indexing set, then

\[
\det (X Y) = (\det X) (\det Y).\tag{8}
\]

For \(X\) to be invertible, it is necessary and sufficient that \(\det X\) be an invertible element of \(\mathbf{A}\) , and then

\[
\det (X ^ {- 1}) = (\det X) ^ {- 1}.\tag{9}
\]

This follows immediately from no. 1, Proposition 1 and no. 2, Theorem 1.

COROLLARY. Two similar square matrices have equal determinants.

If \(P\) is an invertible square matrix, then \(\det(PXP^{-1}) = \det X\) by (8) and (9).

PROPOSITION 6. For the columns of a square matrix X of finite order to be linearly independent, it is necessary and sufficient that det X be not a divisor of zero in A.

This follows from no. 2, Proposition 3.

